\subsection{A multiplicity bound for the quadratic character branch}
\label{sec:quadratic-character-bound}
We now remove the multiplicity restriction from another entire rational
algebra type. The argument uses only two central twist moments, the
orthogonality of rational characters, and the unrestricted exclusion of
$R_n$.

\begin{lemma}[A codegree moment bound]\label{lem:codegree-moment-bound}
Let $R$ be a self-dual rank-four fusion ring with a pseudo-unitary
categorification. Let $d_0=1,d_1,d_2,d_3$ be its Frobenius--Perron
dimensions, $D=\sum_i d_i^2$, and let $f_0,\ldots,f_3$ be its formal
codegrees, counted with multiplicity. Put
\[
 w=\sum_j f_j^{-2},\qquad S=\sum_i d_i.
\]
Then
\begin{equation}\label{eq:codegree-moment-bound}
 D(1-3w)+(1-w)S+4w-\frac2D\geq0,
 \qquad S\leq1+\sqrt{3(D-1)}.
\end{equation}
\end{lemma}
\begin{proof}
Apply the global spherical codegree bound of
\cite[Proposition~\ref{B-prop:spherical-codegree-moment}]{Finite} to the
canonical spherical structure. Its positive simple dimensions make the
absolute-value sum equal to $S$. The second assertion is Cauchy--Schwarz
applied to $d_1,d_2,d_3$.
\end{proof}

\begin{corollary}[A common dimension cutoff]\label{cor:dimension-196}
Under the hypotheses of Lemma~\ref{lem:codegree-moment-bound}, if
$w\geq3/8-2/D$, then $D<196$ and the multiplicity is at most thirteen.
\end{corollary}
\begin{proof}
Suppose $D\geq196$. The left side of
\eqref{eq:codegree-moment-bound} decreases as a function of $w$, since
its coefficient is $-3D-S+4<0$. Thus
\begin{align*}
 0&\leq D(1-3w)+(1-w)S+4w-2/D\\
  &\leq-\frac D8+\frac{15}2+
         \left(\frac58+\frac2D\right)S-\frac{10}D\\
  &<-\frac D8+\frac{49}6+\frac76\sqrt D\\
  &=-\frac{(\sqrt D-14)(3\sqrt D+14)}{24}\leq0.
\end{align*}
The strict estimate uses $5/8+2/D<2/3$ and
$S\leq1+\sqrt{3(D-1)}<1+(7/4)\sqrt D$.
This contradiction proves $D<196$.
Every entry of a real symmetric nonnegative fusion matrix is at most
its operator norm, its Frobenius--Perron eigenvalue $d_i$.
Since $d_i^2<D<196$, every integral structure constant is at most thirteen.
\end{proof}

\begin{lemma}\label{lem:two-rational-codegrees}
Suppose a self-dual rank-four fusion ring has two rational characters with
formal codegrees $m,n$. Then $m,n$ are integers at least two. The pairs
$(3,4)$ and $(3,5)$ are impossible. If $m=n=3$, then the ring is $R_k$ for an integer $k\geq1$.
\end{lemma}
\begin{proof}
The character values are rational algebraic integers, hence integers, and
the first value of every character is one. Its formal codegree is the sum
of the four squared values. Distinct character vectors are orthogonal
because the fusion matrices are commuting real symmetric matrices.
A rational character of codegree three therefore has nontrivial entries
$0,\pm1,\pm1$. A character of codegree four has nontrivial entries
$\pm1,\pm1,\pm1$, while a character of codegree five has entries
$\pm2,0,0$, up to permutation. Its scalar product with a codegree-three
character is odd in both cases, and cannot vanish. This excludes the two
pairs. A codegree equal to one would give a vector $(1,0,0,0)$, which
cannot be orthogonal to the positive Frobenius--Perron character; thus
$m,n\geq2$.

Suppose now $m=n=3$. Orthogonality to the positive dimension vector,
whose entries are all at least one, forces the two nonzero nontrivial
entries of each rational character to have opposite signs. Their zero
coordinates are distinct: otherwise their scalar product would again
be odd. Orthogonality between the two characters then forces a chain of
two differences of Frobenius--Perron dimensions equal to one. After
ordering the basis along this chain, the characters and dimension vector
are
\[
 (1,0,1,-1),\qquad(1,1,-1,0),\qquad
 (1,\delta-1,\delta,\delta+1).
\]
Apply the two rational characters to the six products of the three
nontrivial basis elements, using the symmetric parameter tuple
$(a,b,c,d,e,f,g,h,i,j)$. The first character gives
\[
 c=b+1,\quad e=d,\quad f=e,\quad h=g,\quad i=h+1,\quad j=i;
\]
the second gives
\[
 a=b,\quad d=b+1,\quad e=c,\quad g=d,\quad h=e,\quad i=f+1.
\]
Set $k=d=b+1\geq1$. These linear equations give exactly
\[
 (a,b,c,d,e,f,g,h,i,j)
 =(k-1,k-1,k,k,k,k,k,k,k+1,k+1),
\]
which is the defining multiplication of $R_k$.
\end{proof}

\begin{theorem}\label{thm:quadratic-character-bound}
Let $R$ be a self-dual rank-four fusion ring such that
\[
 R\otimes_{\mathbb Z}\mathbb Q\cong\mathbb Q\times\mathbb Q\times K
\]
for a real quadratic field $K$. If $R$ admits a pivotal categorification,
then its multiplicity is at most thirteen. Consequently the bounded classification certificates through
multiplicity thirteen give the unrestricted pivotal classification
for this rational algebra type.
\end{theorem}
\begin{proof}
Proposition~\ref{prop:rank4-galois-exception} gives a pseudo-unitary Galois
conjugate or an already pseudo-unitary category. We may therefore apply
Lemma~\ref{lem:codegree-moment-bound} to the original Frobenius--Perron
dimensions.

First suppose the Frobenius--Perron character is rational. All its
values are integers. The formal-codegree field condition then makes all
four formal codegrees positive integers. Since their reciprocals sum to
one, their maximum is at most $42$. Here is the elementary four-term
unit-fraction bound: after ordering the denominators $a\leq b\leq c\leq d$,
one has $2\leq a\leq4$. If $a=4$, all four are $4$. If $a=3$,
then $b=3$ or $4$, and the equation for $c,d$ gives $d\leq12$.
If $a=2$, then $3\leq b\leq6$; for $b=3,4,5,6$ the maximum possible
$d$ is respectively $42,20,10,6$. These follow directly from
$1/c+1/d=1-1/a-1/b$, with $c\leq d$ and a positive right side.
As $D$ is the largest formal codegree, $D\leq42$.

Suppose now that the Frobenius--Perron character is quadratic. Its two
characters have codegrees $D,D'$, where $D'$ is the Galois conjugate
of $D$ and is allowed to equal $D$. The
remaining two codegrees $m,n$ come
from the rational characters. Put $u=1/m$, $v=1/n$ and $c=1-u-v$.
The reciprocal identity gives $1/D'=c-1/D$, and therefore
\[
 w=u^2+v^2+(c-1/D)^2+D^{-2}
   =u^2+v^2+c^2-\frac{2c}{D}+\frac2{D^2}.
\]
We have $0<c<1$. If $m=n=3$, Lemma~\ref{lem:two-rational-codegrees}
identifies the ring as $R_k$. The unrestricted quadratic-family theorem
excludes every $k\geq2$; for $k=1$ one has
$D=(39+9\sqrt{13})/2<196$.

For every other possible pair, we claim
\[
 u^2+v^2+(1-u-v)^2\geq\frac38.
\]
If one denominator is two, complete the square in the other reciprocal:
\[
 \frac14+v^2+(\tfrac12-v)^2
 =\frac38+2(v-\tfrac14)^2.
\]
If both denominators are at least four, the expression decreases in each
of $u,v$ on $0<u,v\leq1/4$, so its minimum occurs at $u=v=1/4$ and is
$3/8$. The only remaining case has one denominator three. The other is
neither three, four nor five, by Lemma~\ref{lem:two-rational-codegrees},
and the case two was already covered. Thus $v\leq1/6$, and
\[
 \frac19+v^2+(\tfrac23-v)^2\geq\frac7{18}>\frac38.
\]
This proves the claim and gives $w\geq3/8-2/D$.

Corollary~\ref{cor:dimension-196} now gives $D<196$, completing all cases.

Finally, each fusion matrix $N_i$ is real symmetric and nonnegative, with
spectral radius $d_i$. Hence each of its entries is at most its operator
norm $d_i$. Since $d_i^2<D<196$, every fusion coefficient is less than
fourteen, and its integrality gives the claimed multiplicity bound.
\end{proof}

\begin{corollary}\label{cor:quadratic-character-classification}
The self-dual rank-four fusion rings with rational algebra
$\mathbb Q\times\mathbb Q\times K$, for a real quadratic field $K$,
which admit pivotal categorifications are exactly the four rings of
\[
 \operatorname{Rep}(D_{10}),\qquad
 \operatorname{PSU}(2)_6,\qquad 2D2,\qquad\mathcal H_1,
\]
where $D_{10}$ is the dihedral group of order ten. Each has a unitary
categorification. If the ring has no nontrivial invertible basis object,
only the ring of $\mathcal H_1$ occurs.
\end{corollary}
\begin{proof}
Theorem~\ref{thm:quadratic-character-bound} reduces the assertion to the
exact classification certificates through multiplicity thirteen.
The exact model verifier reads the full tensors in
\path{data/models.json}. For every self-dual model it constructs
$A=N_X+2N_Y+4N_Z$, checks that its characteristic polynomial is
squarefree, and factors it over $\mathbb Q$. Thus $A$ generates
the rational fusion algebra, whose factor degrees identify its
algebra type. The routine
\path{companion_software/verify_unbounded_quadratic_branch.py}
verifies that precisely the four stated models have two rational
factors and one quadratic factor. This computation uses the model
tensors themselves, independently of their census indices.
The first three stated rings have a nontrivial invertible basis
object; $\mathcal H_1$ does not. Their unitary realizations are
those cited in the positive-model table.
\end{proof}

\begin{corollary}[The integral self-dual case]\label{cor:integral-selfdual-rank4}
Every complex-categorifiable self-dual rank-four fusion ring with rational
Frobenius--Perron character has $\operatorname{FPdim}(R)\leq42$ and
multiplicity at most six. The only such rings are those of
$\operatorname{Vec}_{C_2\times C_2}$ and $\operatorname{Rep}(D_{10})$.
\end{corollary}
\begin{proof}
A rational Frobenius--Perron character is integral. Any categorification is
therefore pseudo-unitary, so the rational-character argument in the proof
of Theorem~\ref{thm:quadratic-character-bound} applies without a pivotal
hypothesis: the four formal codegrees are positive integers with reciprocal
sum one, and their maximum $D$ is at most $42$. Every fusion coefficient
is at most $\max_i d_i<\sqrt{42}<7$. The bounded-classification
certificates through multiplicity six leave exactly the two stated
self-dual integral rings, both with standard unitary realizations.
\end{proof}
